% ===== 原文 pp.160–171（PDF p.175–186）=====
% 第 4 章：§4.4 末（Lemma 4.25 起）→ §4.5 → §4.6 → §4.7
% 章标题由 c04a 写，本块不写；本块写 §4.5、§4.6、§4.7 标题。
% 本块无编号公式（tag 区间为空）；文中 (4.8)、(4.15)–(4.18) 均为对 c04b 中编号公式的引用，纯文本书写。
% 边界约定（Lead 裁决）：Theorem 4.24 的陈述由 c04b 写完并闭合，c04c 从 p.160 “Before proving this theorem” 起续写证明。

% [原文 p.160]
在证明本定理之前，我们需要几个引理。

\begin{lemma}
设 $X$ 是环面簇，$W\subseteq X$ 是真闭子集，且没有不可约分支含于 $\partial X$ 中。则存在环面爆破 $\phi:\widetilde X\to X$，使得 $W$ 在 $\widetilde X$ 中的真变换 $\widetilde W$ 不含任何零维分层。
\end{lemma}

\begin{proof}
任何环面爆破都由定义 $X$ 的扇 $\Sigma$ 的某个加细给出。注意：若我们借助 $\Sigma$ 的某个加细得到所需结果，那么 $\Sigma$ 的进一步加细仍然奏效。因此，我们只需对单个锥 $\sigma$ 解决问题：因为若对 $\Sigma$ 中每个锥 $\sigma$ 都有加细 $\Sigma_\sigma$ 在其对应的开仿射子集上奏效，我们就可以选取 $\Sigma$ 的一个加细，使它在每个 $\sigma\in\Sigma$ 上都限制为 $\Sigma_\sigma$ 的加细。
于是可以假设：若 $X$ 是 $n$ 维的，则 $\sigma$ 是 $n$ 维锥且 $X=X_\sigma$（若 $\dim\sigma<n$，则 $X_\sigma$ 没有零维分层）。此外，我们可以用一个包含 $W$、且不含 $X_\sigma$ 的任何环面分层的超曲面替换 $W$。
假设 $W$ 由一个方程
\[
f=\sum_{p\in\sigma^\vee\cap N}a_pz^p\in k[\sigma^\vee\cap N]
\]
定义，并令
\[
\Delta_f=\operatorname{Conv}\Bigl(\bigcup_{p:\ a_p\neq0}\bigl(p+(\sigma^\vee\cap N)\bigr)\Bigr)\subseteq N_\R .
\]
这是 $f$ 的 Newton 多面体。设 $\widetilde\Sigma$ 是 $\Delta_f$ 的法扇。它显然是 $\Sigma$ 的加细，从而定义 $X$ 的一个爆破 $\phi:\widetilde X\to X$。我们需要证明 $\phi$ 具有所需性质。
考虑 $\widetilde\Sigma$ 中的一个 $n$ 维锥 $\tau$；它将是 $\Delta_f$ 某个顶点 $v$ 对应的法锥 $N_{\Delta_f}(v)$。特别地，$\tau^\vee$ 是切锥 $T_v(\Delta_f)$，故 $\Delta_f\subseteq v+\tau^\vee$ 且 $a_v\neq0$。现在 $\phi:X_\tau\to X_\sigma$ 由锥的包含 $\sigma^\vee\to\tau^\vee$ 诱导，而 $f\circ\phi$ 中每个单项式都被 $z^v$ 整除。因此 $W$ 在 $X_\tau$ 中的真变换 $\widetilde W$ 由方程
\[
\sum_{p\in\sigma^\vee\cap N}a_pz^{p-v}=0
\]
定义，且该多项式有非零常数项。由于所有非常数单项式都在 $X_\tau$ 的零维分层上消失，$\widetilde W$ 不含零维分层，正如所欲。
% [Tutor Note: 引理 4.25 的证明用 Newton 多面体 $\Delta_f$ 的法扇给出 $\widetilde\Sigma$；核对锚点 p.160 的 $\Delta_f=\operatorname{Conv}\bigl(\bigcup_{p:a_p\neq0}(p+(\sigma^\vee\cap N))\bigr)$ 与真变换方程 $\sum a_pz^{p-v}=0$（$z^v$ 为所除单项式）。另：第 4 章原文用黑板体 $k$ 表示基域（本译用 \texttt{\textbackslash Bbbk}），$k^\dagger$ 表示相应的对数概形底。]
\end{proof}

% [原文 p.161]
我们需要它的一个稍强版本：

\begin{lemma}
设 $X$ 是环面簇，$W\subseteq X$ 是真闭子集，没有不可约分支含于 $\partial X$ 中。假设 $W$ 在 $X$ 中的余维数 $>c$。则存在环面爆破 $\phi:\widetilde X\to X$，使得 $W$ 在 $\widetilde X$ 中的真变换 $\widetilde W$ 与任何维数 $\le c$ 的环面分层不相交。
\end{lemma}

\begin{proof}
对 $c$ 归纳：我们可以通过作若干次环面爆破，假设 $W$ 与维数小于 $c$ 的环面分层不相交。于是只需证明：可以找到 $\widetilde X$，使 $\widetilde W$ 与任何 $c$ 维环面轨道不相交。现在 $X$ 中一个 $c$ 维环面轨道含于某个仿射开子集 $X_\tau$ 之中，其中 $\tau\in\Sigma$ 且 $\operatorname{codim}\tau=c$。我们可以如引理 4.25 的证明那样论证：只需处理 $X=X_\tau$ 的情形。
于是假设 $\Sigma$ 由 $\tau$ 的面组成，并令 $\overline\Sigma$ 为同一个扇 $\Sigma$，但看作更小子空间 $\R\tau$ 中的扇。任意取定一个线性投影 $M\to(\R\tau)\cap M$，都定义了从 $\Sigma$ 到 $\overline\Sigma$ 的扇映射，从而给出映射 $\psi:X=X_\Sigma\to X_{\overline\Sigma}=:Y$。注意 $\dim Y=\dim X-c$，且 $X$ 中唯一的 $c$ 维分层被映到 $Y$ 中唯一的零维分层。又 $\dim\psi(W)\le\dim W<\dim X-c=\dim Y$，故 $\psi(W)\subseteq Y$ 是真子集，且没有不可约分支含于 $\partial Y$ 中。于是由引理 4.25，存在 $\overline\Sigma$ 的一个细分，给出环面爆破 $\widetilde Y\to Y$，使得 $\psi(W)$ 在 $\widetilde Y$ 中的真变换与 $\widetilde Y$ 的零维环面分层不相交。但 $\overline\Sigma$ 的这个细分当然也给出 $\Sigma$ 的一个细分，从而给出爆破 $\widetilde X\to X$。由于 $W$ 的真变换 $\widetilde W$ 映到 $\psi(W)$ 在 $\widetilde Y$ 中的真变换，$\widetilde W$ 现在与 $\widetilde X$ 中的 $c$ 维环面分层不相交。
\end{proof}

% [Tutor Note: 引理 4.26 的证明有两处易错：(i) 归纳假设是“$W$ 与维数小于 $c$ 的环面分层不相交”，而待处理的是 $c$ 维环面轨道；(ii) $\overline\Sigma$ 是同一扇 $\Sigma$ 在子空间 $\R\tau$ 中的像，$M\to(\R\tau)\cap M$ 的任意线性投影都给出扇映射 $\psi:X\to Y$，且 $\dim Y=\dim X-c$。核对锚点：p.161 的“$\operatorname{codim}\tau=c$”与“$X$ 中唯一的 $c$ 维分层映到 $Y$ 中唯一的零维分层”。]
\begin{proof}[定理 4.24 的证明]
\para{第 1 步。} 我们将证明：在作适当的基变换之后，存在一个图表 (4.17)，使得 $\psi^*(t)$ 是 $R$ 的单值化参数，且对 $\widetilde X_0$ 的每个不可约分支 $\widetilde D_v$，映射 $f^{-1}(\widetilde D_v)\to\widetilde D_v$ 是环面横截稳定映射。
设 $\widetilde C^*\to C^*$ 是 $C^*$ 的正规化。设 $y^*_1,\dots,y^*_p$ 是导子轨迹的点，即 $\widetilde C^*$ 中位于 $C^*$ 的二重点之上的点。这些点未必定义在域 $L$ 上，但定义在某个有限扩张 $L'$ 上，而 $L'$ 必然是 $R[t]/(t^d-u)$ 对某个 $d$ 的分式域。把 $L$ 换成 $L'$，我们可以假设这些点定义在 $L$ 上，于是 $C^*$ 由 $\widetilde C^*$ 上若干对定义在 $L$ 上的点粘合得到。在 $\widetilde C^*$ 的一个给定分支 $D^*$ 上，我们有某个集合 $\{x^*_{ik}\}$，其中的点映到 $C^*$ 的标记点，此外还有一些补充点 $\{y^*_{jk}\}$。于是只需证明：存在一个 $\widetilde X$，它对 $\widetilde C^*$ 中出现的每条（定义在 $\operatorname{Spec}L$ 上的）曲线 $(D^*,x^*_{i_1},\dots,x^*_{i_k},y^*_{j_1},\dots,y^*_{j_\ell})$ 都奏效，即得到环面横截稳定映射 $(D,x_1,\dots,x_{ik},y_{j1},\dots,y_{j\ell})\to\widetilde X$。然后我们可以把所得这些曲线沿按 $y$ 标号的成对截面粘合起来，其方式与 $C^*$ 由按 $y^*$ 标号的成对点粘合得到的方式相同。这样便化归到 $C^*$ 几何不可约的情形。事实上，下文中我们将对每个不可约分支 $C^*$ 构造一个加细 $\widetilde{\mathcal P}$；然后取一个公共加细，按构造它不会破坏每个不可约分支所需的性质。我们还可以通过基变换 $\mathbb{A}^1_k\to\mathbb{A}^1_k$、$t\mapsto t^e$，总假设这些加细只有整顶点。下文中我们将自动这样做，不再说明。

% [原文 p.162]
现在设 $W$ 是 $C^*\to X$ 的像的闭包。则 $W$ 是维数至多为 $2$ 的闭子集；又由于复合 $C^*\xrightarrow{f^*}X\xrightarrow{\pi}\mathbb{A}^1_k$ 是支配的（它同时也是复合 $C^*\to\operatorname{Spec}L\to\mathbb{A}^1_k$），$W$ 的维数至少为 $1$。现在有两种情形。
\para{情形 1。} $\dim W=1$。注意这只能发生在 $f^*$ 在 $C^*$ 上为常值映射的情形，此时 $W$ 就是它的泛点 $f(C^*)$ 的闭包。现在，由于 $\operatorname{codim}W=2$，我们可以用引理 4.26 加细 $\widetilde\Sigma$ 得到 $\widetilde X$，并假设 $W$ 与 $\widetilde X$ 的任何一维环面分层都不相交。此外，$\widetilde X$ 可以选得使 $\widetilde X\setminus\widetilde X_0\cong X\setminus X_0$，因为 $W$ 与 $\partial(X\setminus X_0)$ 不相交：否则 $f^*:C^*\to X\setminus X_0$ 就不是环面横截的。于是我们可以假设 $\widetilde X$ 来自 $\mathcal P$ 的某个加细 $\widetilde{\mathcal P}$。
% [Tutor Note: p.162 情形 1 的“$f^*$ 在 $C^*$ 上为常值映射”并不与 $\dim W=1$ 矛盾：常值映射的像是 $W$ 的泛点（$\pi\circ f^*$ 经 $\operatorname{Spec}L\to\mathbb{A}^1_k$ 分解后仍可支配），$W$ 是该泛点的闭包，故可以是曲线。核对锚点：同页“$\operatorname{codim}W=2$”，以及 p.163 开头由稳定约化得到的 $f:(C,x_1,\dots,x_k)\to\widetilde X$。]
现在注意，由于 $\pi$ 是紧合的，投影 $W\to\mathbb{A}^1_k$ 也是紧合的，并且这个投影同样是支配的。又由于 $f$ 在 $C^*$ 上为常值，映射 $f^*:C^*\to W$ 经过 $\operatorname{Spec}L$ 分解。我们此时诉诸稳定曲线的稳定约化（命题 2.4 的一个特殊情形，应用于目标空间 $X$ 是一个点的情形），再次在可能作基变换之后，得到稳定曲线族 $(C,x_1,\dots,x_s)\to\operatorname{Spec}R$。对 $W\to\mathbb{A}^1_k$ 应用紧合性的赋值判别法，我们得到交换图表
\[
\begin{array}{ccc}
\operatorname{Spec}R & \longrightarrow & W\\
\big\downarrow & & \big\downarrow\\
\operatorname{Spec}L & \longrightarrow & \mathbb{A}^1_k
\end{array}
\]
而复合 $C\to\operatorname{Spec}R\to W\hookrightarrow\widetilde X$ 现在就给出所需的映射 $f:(C,x_1,\dots,x_s)\to\widetilde X$。由于 $W$ 与任何一维环面分层都不相交，对 $\widetilde X_0$ 中与 $W$ 相交的每个不可约分支 $\widetilde D_v$，映射 $f^{-1}(\widetilde D_v)\to\widetilde D_v$ 是环面横截稳定映射；而对 $\widetilde X_0$ 的其它不可约分支，它是空的。
\para{情形 2。} $\dim W=2$。此时由引理 4.26，我们可以细分 $\Sigma$ 得到 $\widetilde X$，使得 $W$ 的真变换与 $\widetilde X$ 的零维分层不相交；同样可以做到 $\widetilde X\setminus\widetilde X_0=X\setminus X_0$。于是我们假设这来自 $\mathcal P$ 的某个加细 $\widetilde{\mathcal P}$。
设 $\tau\in\widetilde{\mathcal P}$ 是一条边，$D_\tau\subseteq\widetilde X$ 是相应的一维分层，$X_\tau:=X_{C(\tau)}$ 是 $\widetilde X$ 中相应的开仿射子集。则 $X_\tau\cong\mathbb{G}_m\times V_e$，其中
\[
V_e=\operatorname{Spec}k[x,y,t]/(xy-t^e),
\]
而 $e$ 是 $\tau$ 的仿射长度。设 $C^*_\tau=f^{-1}(X_\tau)\subseteq C^*$，并把 $C^*_\tau\xrightarrow{f}X_\tau$ 与投影 $X_\tau\to V_e$ 复合，得到映射 $h:C^*_\tau\to V_e$。存在开子集 $U\subseteq C^*_\tau$，使 $h$ 在 $U$ 上是 étale 的；若 $h$ 支配，则 $U$ 非空（见 [83], Proposition 3.8）。设 $Z_\tau\subseteq V_e$ 是包含 $C^*_\tau\setminus U$ 的像与 $C^*_\tau$ 的标记点的像的最小闭子集。由于 $C^*_\tau$ 的像与 $\partial V_e$（由 $t=0$ 给出）不相交，$Z_\tau$ 的不可约分支没有含于 $\partial V_e$ 中的。于是我们可以应用引理 4.26，爆破 $V_e$，使 $Z_\tau$ 的真变换与 $V_e$ 的零维分层不相交。这对应于边 $\tau$ 的一个细分，因此我们可以选取 $\widetilde{\mathcal P}$ 的进一步细分，使其在 $\tau$ 上诱导出这个细分。这样做之后，我们可以假设 $Z_\tau$ 与 $V_e$ 的零维分层不相交。我们对每个满足 $\dim\tau=1$ 且 $W\cap D_\tau\neq\varnothing$ 的 $\tau\in\widetilde{\mathcal P}$ 都这样做。这就在此情形下给出了 $\widetilde X$。

% [原文 p.163]
为什么这有效？首先，由稳定映射的稳定约化定理（命题 2.4），在可能再进行一次基变换后，我们得到一个在 $\operatorname{Spec}R$ 上延拓 $f:(C^*,x^*_1,\dots,x^*_k)\to\widetilde X$ 的带标记稳定映射 $f:(C,x_1,\dots,x_k)\to\widetilde X$。现在这个延拓后的 $f$ 的像含于 $W$ 在 $\widetilde X$ 中的真变换 $\widetilde W$ 之中，而按构造 $\widetilde W$ 避开 $\widetilde X$ 的零维分层。特别地，$f(C_0)$ 的任何不可约分支都不可能含于 $\widetilde X$ 的一维环面分层中，因为那样它也会含有 $\widetilde X$ 的一个零维分层。
这还不够：我们仍需证明 $C_0$ 的任何不可约分支都不被 $f$ 收缩成 $\widetilde X$ 的一维分层中的一个点。所以假设确有这样一个不可约分支，被收缩到由 $\tau\in\widetilde{\mathcal P}$ 标号的 $\widetilde X$ 的一维分层中。这如前定义了 $X_\tau$、$C_\tau=f^{-1}(X_\tau)\subseteq C$，以及 $X_\tau\cong\mathbb{G}_m\times V_e$（对某个 $e$）。我们有如上定义的集合 $Z_\tau\subseteq V_e$，它与 $V_e$ 的奇点不相交。注意，若投影 $h:C^*_\tau\to V_e$ 不支配，则 $Z_\tau$ 包含 $C^*_\tau$ 的像，从而包含 $\widetilde W\cap X_\tau$ 在投影到 $V_e$ 下的像。特别地，此时 $\widetilde W$ 与 $X_\tau$ 的一维分层不相交。于是我们可以假设 $h$ 是支配的。设 $\widetilde Z_\tau\subseteq X_\tau$ 是 $Z_\tau$ 在投影 $X_\tau\to V_e$ 下的原像，并令 $Z'_\tau=f^{-1}(\widetilde Z_\tau)$。映射 $f:C_\tau\setminus Z'_\tau\to\widetilde X$ 分解为
\[
C_\tau\setminus Z'_\tau\xrightarrow{\ f'\ }\operatorname{Spec}R\times_{\mathbb{A}^1_k}(X_\tau\setminus\widetilde Z_\tau)\to\widetilde X,
\]
其中 $f'$ 是紧合的。我们知道，$f'$ 除在某个有限（可能为空）的子集
\[
T\subseteq\operatorname{Spec}R\times_{\mathbb{A}^1_k}(X_\tau\setminus\widetilde Z_\tau)
\]
之上以外是有限的，因为 $f'$ 紧合，且除在 $T$ 之上以外是拟有限的（这要用到 $h$ 支配）。现在考虑 $f'$ 的 Stein 分解
\[
C_\tau\setminus Z'_\tau\xrightarrow{\ f''\ }Y_\tau\xrightarrow{\ g\ }\operatorname{Spec}R\times_{\mathbb{A}^1_k}(X_\tau\setminus\widetilde Z_\tau)\to\widetilde X,
\]
其中 $f''$ 在 $(f')^{-1}(T)$ 之外是同构。于是 $Y_\tau\to\widetilde X$ 与 $f|_{C\setminus f^{-1}(T)}$ 粘合，给出映射 $g':C'\to\widetilde X$。这个映射由复合 $x_i:\operatorname{Spec}R\to C\to C'$ 标记，因为按构造 $C_\tau\setminus Z'_\tau$ 不含任何标记点。我现在断言：这个新映射 $g':(C',x_1,\dots,x_k)\to\widetilde X$ 是稳定映射；这与 $f:(C,x_1,\dots,x_k)\to\widetilde X$ 的稳定性矛盾，除非 $T$ 为空。
为此，我们需要说明 $C'_0$ 至多只有二重点。考虑复合映射 $h:Y_\tau\to\operatorname{Spec}R\times_{\mathbb{A}^1}(V_e\setminus Z_\tau)$。按构造，这个映射是有限的，并且在由 $t=0$ 给出的轨迹之外是 étale 的。另一方面，若它在 $t=0$ 的某个分支上分歧，则 $Y_\tau\to\operatorname{Spec}R$ 在 $p$ 上的纤维会有非约化分支，而这是不可能的，因为 $f:C\to\widetilde X$ 是稳定映射。因此 $h$ 除在 $\operatorname{Spec}R\times_{\mathbb{A}^1}(V_e\setminus Z_\tau)$ 的一个有限子集上之外是 étale 的。事实上，由分歧轨迹的纯性（例如见 [1], X, 3.1），这个映射可能不是 étale 的唯一可能点，就是这个 $\operatorname{Spec}R\times_{\mathbb{A}^1}(V_e\setminus Z_\tau)$ 的奇点。然而，étale 局部地，在该奇点的补集上 étale 的有限映射都形如典范覆叠 $V_{e'}\to V_e$（其中 $e'|e$）。于是可以看出，$Y_\tau$ 在 $p$ 上的纤维至多只有二重点。
% [Tutor Note: p.163 的 Stein 分解链 $C_\tau\setminus Z'_\tau\xrightarrow{f''}Y_\tau\xrightarrow{g}\operatorname{Spec}R\times_{\mathbb{A}^1}(X_\tau\setminus\widetilde Z_\tau)$ 中，$f'$ 紧合且除有限集 $T$ 之上外拟有限，$f''$ 在 $(f')^{-1}(T)$ 之外是同构；反证时把 $g'$ 的稳定性与 $f$ 的稳定性对比。核对锚点：“$C'_0$ 至多二重点”的论证用到分歧轨迹纯性（[1], X, 3.1）与典范覆叠 $V_{e'}\to V_e$（$e'|e$）。]
这就完成了情形 2 的论证。
现在考虑我们已经构造出的图表 (4.17) 中的映射 $\psi:\operatorname{Spec}R\to\mathbb{A}^1_k$。假设对某个 $d\ge1$ 有 $\psi^*(t)\in\mathfrak m_R^d\setminus\mathfrak m_R^{d-1}$，于是 $\psi^*(t)=u^d$，其中 $u$ 是单值化参数。若 $d=1$，我们就完成了。否则作 $d$ 次基扩张 $\mathbb{A}^1_k\to\mathbb{A}^1_k$：此时 $R\times_{\mathbb{A}^1}\mathbb{A}^1=R[t]/(t^d-u^d)$ 由 $d$ 个不可约分支组成，

% [原文 p.164]
每个都同构于 $\operatorname{Spec}R$。然后我们可以用其中任一不可约分支，得到满足 $\psi^*(t)=u$ 的图表 (4.17)。
\para{第 2 步。} 得到图表 (4.17) 之后，我们现在要证明 $f$ 在 $C_0$ 上的限制 $f_0$ 是环面横截前对数曲线。由第 1 步，只需证明命题 4.9 的结论 (1) 与 (2) 成立。
设 $p\in C_0$ 是映到 $\widetilde X_0$ 的奇点轨迹的闭点。考虑完全局部 $k[[t]]$-代数的诱导同态
\[
f^*_0:\widehat{\mathcal O}_{\widetilde X,f_0(p)}\to\widehat{\mathcal O}_{C,p}.
\]
现在 $p$ 或者是 $C_0$ 的光滑点，或者是 $C_0$ 的二重点。在前一情形，映射 $C\to\operatorname{Spec}R$ 本身在 $p$ 处光滑，从而由 $f^{-1}(\partial\widetilde X)\subseteq C$ 在 $C$ 上诱导的对数结构在 $p$ 处是对数光滑的。把这个对数结构限制到 $C_0$，得到在 $p$ 处对数光滑的曲线 $C^\dagger_0\to\operatorname{Spec}k^\dagger$，连同 $\operatorname{Spec}k^\dagger$ 上的对数态射 $C^\dagger_0\to\widetilde X^\dagger_0$。但由于 $f(p)$ 位于 $\widetilde X_0$ 的奇点轨迹中，这与命题 4.9 矛盾（那里的论证是纯局部的）。因此我们可以假设 $p$ 是 $C_0$ 的二重点。注意
\[
\widehat{\mathcal O}_{C_0,p}=k[[x,y]]/(xy),
\]
由此推出
\[
\widehat{\mathcal O}_{C,p}=k[[x,y,t]]/(xy-ft^e)
\]
对某个 $e>0$ 与某个 $f\in k[[x,y,t]]$ 成立。容易验证：这个 $k[[t]]$-代数于是同构于
\[
k[[x,y,t]]/(xy-\lambda t^e),
\]
其中 $\lambda$ 为 $0$ 或 $1$，$e>0$。
若 $\lambda=0$，则 $C$ 在 $p$ 的邻域内（在 étale 拓扑下）是可约的。把 $f$ 限制到其中一个分支，便如在 $p$ 是 $C$ 的光滑点的情形中一样得到矛盾。
若 $\lambda=1$，则再次取由 $f^{-1}(\partial\widetilde X)\subseteq C$ 在 $C$ 上诱导的对数结构，$C^\dagger\to\operatorname{Spec}R^\dagger$ 在 $p$ 的一个邻域内是对数光滑的，于是我们可以应用命题 4.9 得出 $f_0$ 是环面横截前对数曲线。这就完成了定理 4.24 的 (1) 的证明。
\para{第 3 步。} 剩下的只需证明：在定理 4.24 的 (2) 的情形中，$C^\dagger\to\operatorname{Spec}R^\dagger$ 事实上是对数光滑的。事实上，第 2 步的论证已经表明，在映到 $\operatorname{Sing}(X_0)$ 的所有点 $p\in C_0$ 的邻域内情形如此。对于不映到 $\operatorname{Sing}(X_0)\cup\partial X_0$ 的任何光滑点 $p\in C_0$，情形同样如此；而条件“$f^{-1}(\partial(\widetilde X\setminus\widetilde X_0))$ 是 $C\to\operatorname{Spec}R$ 的一族互不相交截面的并”，给出了在映到 $\partial\widetilde X_0$ 的光滑点 $p$ 处的对数光滑性。
最后，与前对数曲线 $f_0$ 相伴的热带曲线 $h:\Gamma\to M_\R$ 是简单曲线这一假设蕴含：若 $v\in\mathcal P$ 是满足 $h^{-1}(v)$ 非空的顶点，则要么 $h^{-1}(v)$ 是 $\Gamma$ 的单个三价顶点，要么 $h^{-1}(v)$ 由 $\Gamma$ 的若干条边内部的点组成。在前一情形，$f^{-1}(D_v)$（因 $C_0$ 有理而有理）由一条单直线组成；在后一情形，$f^{-1}(D_v)$ 由若干互不相交的二价直线组成。特别地，$C_0$ 没有不映到 $\operatorname{Sing}(X_0)$ 的二重点。于是 $C^\dagger_0\to\operatorname{Spec}k^\dagger$ 是对数光滑的，正如所欲。
% [Tutor Note: p.164 第 3 步中原文用 $\operatorname{Sing}(X_0)$，而第 2 步用 $\widetilde X_0$ 的“奇点轨迹”（singular locus of $\widetilde X_0$）。按上下文，第 3 步的 $\operatorname{Sing}(X_0)$ 应理解为 $\operatorname{Sing}(\widetilde X_0)$；照译不改，见 PDF p.179 图像。]
\end{proof}

% [原文 p.165]
\subsection{对数世界 $\to$ 经典世界}
我们回到基本情形：数据 $P_1,\dots,P_s\in M$，一个良好的格多面体分解 $\mathcal P$，以及一般点 $Q_1,\dots,Q_s\in\mathbb{G}(\widetilde M)$，其中 $s=|\Delta|-1$。这给出截面 $\sigma_1,\dots,\sigma_s:\mathbb{A}^1_k\to X$ 与点 $q_1,\dots,q_s\in X_0$，且 $q_i=\sigma_i(0)$。
本节的主要定理是

\begin{theorem}
设 $f_0:(C^\dagger_0,x_1,\dots,x_s)\to X^\dagger_0$ 是亏格零的环面横截对数曲线，且 $f(x_i)=q_i$，其相伴热带曲线为 $h:\Gamma\to M_\R$。假设 $h$ 是简单的。则存在唯一的、$\operatorname{Spec}k[[t]]$ 上的带标记有理曲线 $(C_\infty,x^\infty_1,\dots,x^\infty_s)$，连同交换图表
\[
\begin{array}{ccc}
(C_\infty,x^\infty_1,\dots,x^\infty_s) & \xrightarrow{\ f_\infty\ } & X\\[2pt]
\big\downarrow & & \big\downarrow\\[2pt]
\operatorname{Spec}k[[t]] & \xrightarrow{\ \psi_\infty\ } & \mathbb{A}^1
\end{array}
\]
使得
\begin{enumerate}[label=(\arabic*),leftmargin=2.2em]
\item $\psi_\infty$ 由自然包含 $k[t]\hookrightarrow k[[t]]$ 诱导；
\item $\sigma_i\circ\psi_\infty=f_\infty\circ x^\infty_i$；
\item 若给 $C_\infty$ 配置由 $f^{-1}_\infty(\partial X)\subseteq C_\infty$ 诱导的对数结构，并给 $\operatorname{Spec}k[[t]]$ 配置从 $\mathbb{A}^1$ 拉回的对数结构，则在 $\operatorname{Spec}k^\dagger$ 上诱导的对数态射 $(C_0,x_1,\dots,x_s)\to X_0$ 与给定的 $f_0$ 重合。
\end{enumerate}
\end{theorem}

% [Tutor Note: 定理 4.27 中的 $C_\infty$ 是 $\operatorname{Spec}k[[t]]$ 上的带标记有理曲线；条件 (1) 说明 $\psi_\infty$ 由 $k[t]\hookrightarrow k[[t]]$ 诱导，即 $C_\infty$ 是把 $f_0$ 沿 $\mathbb{A}^1$“形式地”展成的唯一族；条件 (3) 的“诱导对数态射”是把 $C_\infty$ 的对数结构限制到闭点纤维 $C_0$ 上。核对锚点：定理 4.24 的图表 (4.17)/(4.18)。另 p.165 的一般点集写作 $\mathbb{G}(\widetilde M)$（黑板体 G），与 $\mathbb{G}_m$ 不是同一记号。]
\begin{proof}
\para{第 1 步。有限阶形变理论。} 我们将应用 §3.4 的形变理论，特别是定理 3.43。首先注意：定理 3.43 的假设适用于数据
\[
[f_0:C_0/O_0\to X,\mathbf{x}^0],
\]
其中 $f_0:C^\dagger_0\to X^\dagger$ 现在是原来的 $f_0$ 与包含 $X_0\hookrightarrow X$ 的复合，而 $x^0_i=x_i$。当然 $C_0$ 按假设是有理的，且 $f_{0*}:\Theta_{C^\dagger_0/k^\dagger}\to f^*_0\Theta_{X^\dagger/(\mathbb{A}^1)^\dagger}$ 是单射。确实，$f_{0*}$ 在 $C_0$ 的每个泛点处是单射，因为 $\varphi_0$ 在每个泛点的一个邻域内 étale 到其像上。由于 $f_0$ 对数光滑，$\Theta_{C^\dagger_0/k^\dagger}$ 是局部自由的，从而不可能含有挠子层。因此 $\ker f_{0*}=0$。
我们现在将证明：定理 3.43 中的映射 $\Xi$ 事实上是同构。这就表明：对任何满足 $\sigma_i\circ\psi_{k-1}=f_{k-1}\circ x^{k-1}_i$ 的提升
\[
[f_{k-1}:C_{k-1}/O_{k-1}\to X,\mathbf{x}^{k-1}],
\]
都存在唯一的提升
\[
[f_k:C_k/O_k\to X,\mathbf{x}^k]
\]
使得 $\sigma_i\circ\psi_k=f_k\circ x^k_i$。

% [原文 p.166]
首先注意，我们有图表
\[
\begin{array}{ccccccc}
 & & & & & 0 & \\
 & & & & & \big\downarrow & \\
 & 0 & & & & \Theta_{C^\dagger_0/k^\dagger}|_{\{x_1,\dots,x_s\}} & \\
 & \big\downarrow & & & & \big\downarrow & \\
0\to & \Theta_{C^\dagger_0/k^\dagger}(-\mathbf{x}) & \to & f_0^*\Theta_{X^\dagger/(\mathbb{A}^1)^\dagger} & \xrightarrow{\ \ } & \mathcal N_{f_0,\mathbf{x}} & \to 0\\
 & \big\downarrow & & \big\downarrow{\scriptstyle =} & & \big\downarrow & \\
0\to & \Theta_{C^\dagger_0/k^\dagger} & \to & f_0^*\Theta_{X^\dagger/(\mathbb{A}^1)^\dagger} & \xrightarrow{\ \ } & \mathcal N_{f_0} & \to 0\\
 & \big\downarrow & & & & \big\downarrow & \\
 & \Theta_{C^\dagger_0/k^\dagger}|_{\{x_1,\dots,x_s\}} & & & & 0 & \\
 & \big\downarrow & & & & & \\
 & 0 & & & & &
\end{array}
\]
% [Tutor Note: p.166 蛇引理图表的读法：左列竖直正合列 $0\to\Theta_{C^\dagger_0/k^\dagger}(-\mathbf{x})\to\Theta_{C^\dagger_0/k^\dagger}\to\Theta_{C^\dagger_0/k^\dagger}|_{\{x_1,\dots,x_s\}}\to0$，右列竖直正合列 $0\to\Theta_{C^\dagger_0/k^\dagger}|_{\{x_1,\dots,x_s\}}\to\mathcal N_{f_0,\mathbf{x}}\to\mathcal N_{f_0}\to0$，中央 $f_0^*\Theta_{X^\dagger/(\mathbb{A}^1)^\dagger}=f_0^*\Theta_{X^\dagger/(\mathbb{A}^1)^\dagger}$ 为恒等（标“=”）。蛇引理给出 $\mathcal N_{f_0,\mathbf{x}}\cong\Theta_{C^\dagger_0/k^\dagger}|_{\{x_1,\dots,x_s\}}\oplus\mathcal N_{f_0}$。]
由蛇引理，于是 $\mathcal N_{f_0,\mathbf{x}}$ 分解为天楼层 $\Theta_{C^\dagger_0/k^\dagger}|_{\{x_1,\dots,x_s\}}$ 与 $\mathcal N_{f_0}$ 的直和。若把 $\Xi$ 限制到 $\Theta_{C^\dagger_0/k^\dagger}\otimes k(x_i)$ 上，它恰好由
\[
f_{0*}:\Theta_{C^\dagger_0/k^\dagger}\otimes k(x_i)\to T_{X/\mathbb{A}^1,\sigma_i(0)}
\]
给出。令 $T_{C_0/k,x_i}:=\Theta_{C^\dagger_0/k^\dagger}\otimes k(x_i)$，即 $C_0$ 在 $x_i$ 处的 Zariski 切空间。于是要证明 $\Xi$ 是同构，只需证明
\[
\Xi':H^0(C,\mathcal N_{f_0})\to\prod_{i=1}^sT_{X/\mathbb{A}^1,\sigma_i(0)}/f_{0*}(T_{C_0/k,x_i})
\]
是同构。我们现在更明确地计算 $\Xi'$ 的定义域与靶。
按惯例，设
\[
h:\widetilde\Gamma\to M_\R
\]
由 $\Gamma$ 添加顶点得到，使得对点 $y\in\widetilde\Gamma$，$h(y)$ 是 $\mathcal P$ 的顶点当且仅当 $y$ 是 $\widetilde\Gamma$ 的顶点，或 $y$ 含于某条带标记的无界边的内部。又设
\[
h:\widehat\Gamma\to M_\R
\]
由 $\Gamma$ 去掉无界带标号边、并去掉由此产生的任何二价顶点得到。
靶不难描述：由于 $\Theta_{X^\dagger/(\mathbb{A}^1)^\dagger}$ 典范地同构于 $M\otimes\mathcal O_X$，我们有典范同构
\[
T_{X/\mathbb{A}^1,\sigma_i(0)}\cong M\otimes_\Z k .
\]
另一方面，设 $V_i$ 是 $\widetilde\Gamma$ 中作为标记边 $E_{x_i}$ 的边界的顶点。则曲线 $C_{V_i}$ 是一条直线。设 $u_i$ 是 $h$ 的像中与 $V_i$ 相邻的任一无标记边的切向量。则由引理 4.17 可知，$f_0(C_{V_i})$ 是 $\mathbb{G}(Zu_i)\subseteq\mathbb{G}(M)\subseteq\mathbb{G}(\widetilde M)$ 的一条轨道。于是 $f_{0*}T_{C_0/k,x_i}=$

% [原文 p.167]
$(Zu_i)\otimes_\Z k\subseteq M\otimes_\Z k$。所以 $\Xi'$ 的靶是
\[
\prod_{i=1}^s(M/Zu_i)\otimes_\Z k=\prod_{i=1}^s\bigl(M/Zu_{(\partial_-E_i,E_i)}\bigr)\otimes_\Z k,
\]
这里用了命题 4.22 的记号。
接下来描述 $\Xi'$ 的定义域 $H^0(C_0,\mathcal N_{f_0})$。首先，我们把 $\mathcal N_{f_0}$ 限制到 $C_0$ 的每个分支 $C_V$ 上；注意正合列
\[
0\to\Theta_{C^\dagger_0/k^\dagger}\to f_0^*\Theta_{X^\dagger_0/k^\dagger}=M\otimes_\Z\mathcal O_C\to\mathcal N_{f_0}\to0
\]
限制为如下正合列
\[
0\to\Theta_{C^\dagger_0/k^\dagger}|_{C_V}\to M\otimes_\Z\mathcal O_{C_V}\to\mathcal N_{f_0}|_{C_V}\to0 .
\]
由例 3.36 的 (6)，$\Theta_{C^\dagger_0/k^\dagger}|_{C_V}$ 恰好是 $\Theta_{C_V/k}\bigl(-\sum p_i\bigr)$，其中 $\{p_i\}$ 是 $C_V$ 上特殊点的集合，即 $C^\dagger_0$ 上例 3.26 中 (2) 型与 (3) 型的点。这样的点的个数恰好是 $V$ 在 $\widetilde\Gamma$ 中的价数（不计带标记的边）。由于这个价数（不计带标记的边）是 $2$ 或 $3$，我们就把这两种可能分别称为二价情形与三价情形。于是 $\Theta_{C^\dagger_0/k^\dagger}|_{C_V}\cong\mathcal O_{\PP^1}$（二价情形）或 $\mathcal O_{\PP^1}(-1)$（三价情形）。
再次注意：在二价情形，由于 $f_0(C_V)$ 是 $\mathbb{G}(Zu)$ 的一条轨道（其中 $u$ 与 $V$ 相邻的任一条边相切），$H^0(C_V,\Theta_{C^\dagger_0/k^\dagger}|_{C_V})$ 在 $M\otimes_\Z k$ 中的像是 $(Zu)\otimes k$。于是在二价与三价情形分别有
\[
H^0(C_V,\mathcal N_{f_0}|_{C_V})=
\begin{cases}
(M/Zu)\otimes_\Z k,&V\ \text{为二价},\\[2pt]
M\otimes_\Z k,&V\ \text{为三价}.
\end{cases}
\]
为了理解能否把一族截面
\[
(s_V)\in\bigoplus_{V\in\widetilde\Gamma[0]}H^0(C_V,\mathcal N_{f_0}|_{C_V})
\]
粘成 $H^0(C_0,\mathcal N_{f_0})$ 的一个截面，我们将证明：在与 $\widetilde\Gamma$ 的边 $E$ 对应的二重点 $x\in C_0$ 处，$f_{0*}$ 把一维空间 $\Theta_{C^\dagger_0/k^\dagger}\otimes k(x)$ 同构地映到 $\Theta_{X^\dagger/k^\dagger}\otimes k(f_0(x))=M\otimes_\Z k$ 中由 $h(E)$ 的切向量张成的子空间上，而该子空间可描述为 $Zu_{(\partial_-E,E)}\otimes_\Z k\subseteq M\otimes_\Z k$。
确实，为看清这一点，照惯例把 $C$ 在局部写成 $\operatorname{Spec}k[z,w]/(zw)$，并把 $X$ 在 $f(x)$ 的一个邻域内局部写成 $\operatorname{Spec}k[u,v,y^{\pm},t]/(uv-t^\ell)$（对某个 $\ell$，其中 $\ell$ 是 $h(E)$ 的仿射长度）。我们可以假设 $f^*_0u=z^\mu$、$f^*_0v=w^\mu$，其中 $\mu=w(E)$；又 $f^*_0(t)=0$，$f^*_0(y)=g(z,w)$。则 $\Theta_{C^\dagger/k^\dagger}\otimes k(x)$ 的一个局部生成元是 $z\partial_z-w\partial_w$（见例 3.36 的 (6)），而 $\Theta_{X^\dagger/(\mathbb{A}^1)^\dagger}$ 在 $f(x)$ 处的局部生成元是 $u\partial_u-v\partial_v$ 与 $\partial_y$。容易看出，在 $x$ 处有 $f_{0*}(z\partial_z-w\partial_w)=\mu(u\partial_u-v\partial_v)$。作为 $\Theta_{X^\dagger/k^\dagger}\otimes k(f_0(x))=M\otimes_\Z k$ 的元素，$u\partial_u-v\partial_v$ 被看作 $M$ 中与边 $h(E)$ 相切的元素，因为这个导子在 $y$ 上消失。
因此，给定 $(s_V)\in\bigoplus H^0(C_V,\mathcal N_{f_0}|_{C_V})$，$s_V$ 与 $s_{V'}$ 在二重点 $x\in C_V\cap C_{V'}$ 处可以粘合，当且仅当它们作为 $M\otimes_\Z k$ 的元素相差一个 $Zu_{(\partial_-E,E)}$ 中的元素。
% [原文 p.168]
于是映射
\[
\widetilde\Phi:\prod_{V\in\widehat\Gamma[0]}M\otimes_\Z k\times\prod_{V\in\widetilde\Gamma[0]\setminus\widehat\Gamma[0]}\bigl(M/Zu_{(\partial_-E(V),E(V))}\bigr)\otimes_\Z k\to\prod_{E\in\widetilde\Gamma[1]\setminus\widetilde\Gamma[1]_\infty}\bigl(M/Zu_{(\partial_-E,E)}\bigr)\otimes_\Z k
\]
（由 $(s_V)_{V\in\widetilde\Gamma[0]}\mapsto(s_{\partial_+E}-s_{\partial_-E})_{E\in\widetilde\Gamma[1]\setminus\widetilde\Gamma[1]_\infty}$ 给出）的核正是 $H^0(C_0,\mathcal N_{f_0})$。容易验证 $\widetilde\Phi$ 是满射。于是 $\Xi'$ 是同构，当且仅当
\[
\Phi'':\prod_{V\in\widehat\Gamma[0]}(M\otimes_\Z k)\times\prod_{V\in\widetilde\Gamma[0]\setminus\widehat\Gamma[0]}\bigl(M/Zu_{(\partial_-E(V),E(V))}\bigr)\otimes_\Z k\to\prod_{E\in\widetilde\Gamma[1]\setminus\widetilde\Gamma[1]_\infty}\bigl(M/Zu_{(\partial_-E,E)}\bigr)\otimes_\Z k\times\prod_{i=1}^s\bigl(M/Zu_{(\partial_-E_i,E_i)}\bigr)\otimes_\Z k
\]
是同构。这个映射由
\[
(s_V)_{V\in\widetilde\Gamma[0]}\mapsto\Bigl((s_{\partial_+E}-s_{\partial_-E})_{E\in\widetilde\Gamma[1]\setminus\widetilde\Gamma[1]_\infty},\ (s_{V_i})_{1\le i\le s}\Bigr)
\]
给出，其中 $V_i$ 是 $E_{x_i}$ 的顶点。
% [Tutor Note: p.168 中 $\widetilde\Phi$ 与 $\Phi''$ 的第一个乘积的指标集是 $\widehat\Gamma[0]$（$\widehat\Gamma$ 为去掉无界带标号边后所得图），第二个乘积的指标集是 $\widetilde\Gamma[0]\setminus\widehat\Gamma[0]$（无界边的叶顶点），而 $\widetilde\Gamma[1]_\infty$ 为无界边集。$\Phi''$ 与 (4.8) 的映射相差一个与 $k$ 的张量，见 c04b。]
这个映射与 (4.8) 中的映射 $\Phi''$ 在把后者与 $k$ 张量之后一致。我们在命题 4.22 的证明中已经看到，这确实是同构。于是 $\Xi$ 是同构，从而对任何提升 $\mathcal O^\dagger_{k-1}$ 上对数结构的对数结构 $\mathcal O^\dagger_k$ 的选择，都存在唯一的提升
\[
[f_k:C_k/O_k\to X,\mathbf{x}^k]
\]
作为
\[
[f_{k-1}:C_{k-1}/O_{k-1}\to X,\mathbf{x}^{k-1}]
\]
的提升。
\para{第 2 步。取极限。} 现在使用 $\mathcal O_k$ 上由图卡
\[
\mathbb{N}\to k[t]/(t^{k+1}),\qquad n\mapsto t^n
\]
给出的对数结构。由第 1 步，这给出对一切 $k$ 的 $[f_k:C_k/O_k\to X,\mathbf{x}^k]$。忘掉 $C^\dagger_k\to\mathcal O^\dagger_k$ 上的对数结构，$f_k:C_k/O_k\to X$ 可以看作 $X$ 中稳定映射的族，因而给出从 $\mathcal O_k$ 到亏格零、带 $s$ 个标记点、目标为 $X$ 的稳定映射叠的映射。这些映射是相容的，因此给出从 $\mathcal O_\infty=\operatorname{Spec}k[[t]]$ 到稳定映射叠的映射，即给出图表
\[
\begin{array}{ccc}
C_\infty & \xrightarrow{\ f_\infty\ } & X\\[2pt]
\big\downarrow & & \big\downarrow\\[2pt]
\mathcal O_\infty & \longrightarrow & \mathbb{A}^1
\end{array}
\]
其中 $f_\infty$ 是曲线的稳定映射。这就给出了定理 4.27 中所需的图表。
我们还需要证明这个图表满足定理的条件 (3)。我们首先论证：$C^\dagger_\infty\to\mathcal O^\dagger_\infty$（$C_\infty$ 上的对数结构由包含 $f^{-1}_\infty(\partial X)\subseteq C_\infty$ 诱导）事实上是对数光滑的。这只需在 $C_0$ 的二重点与标记点的 étale 邻域内验证。

% [原文 p.169]
在二重点 $q$ 处，如定理 4.24 的证明中第 2 步所观察到的，我们有
\[
\widehat{\mathcal O}_{C_\infty,q}\cong k[[x,y,t]]/(xy-\lambda t^e),
\]
其中 $\lambda\in\{0,1\}$。把例 3.26 的对数光滑性刻画应用于 $C^\dagger_k\to\mathcal O^\dagger_k$（$k$ 充分大），可知事实上 $\lambda=1$，且 $e$ 由条件 $\overline{\mathcal M}_{C_0,\bar q}=S_e$ 决定。于是 $C^\dagger_\infty\to\mathcal O^\dagger_\infty$ 在 $q$ 处对数光滑。至于标记点 $p$ 处的对数光滑性，只需注意 $f^{-1}_\infty(\partial(X\setminus X_0))$ 如定理 4.24 的证明第 3 步那样是一族互不相交截面的并。这成立是因为对 $f^{-1}_k(\partial(X\setminus X_0))$ 而言同样必然成立，而 $f_k:C^\dagger_k\to\mathcal O^\dagger_k$ 是对数光滑的。于是我们看到 $f_\infty$ 在诱导对数结构下是对数光滑的。
现在我们可以把这个对数结构限制到 $C_0$，得到 $\operatorname{Spec}k^\dagger$ 上另一条环面横截对数曲线 $(C'_0)^\dagger\to X^\dagger_0$。它事实上必定与原来的曲线 $C^\dagger_0\to X^\dagger_0$ 重合。确实，由命题 4.23 的证明，$C_0$ 上的对数结构在结点之外被唯一确定，而在结点处例 3.26 的局部描述告诉我们两个对数结构重合。这证明了存在性。
至于唯一性，给定这样的 $f_\infty:C_\infty\to X$，我们得到 $\mathcal O^\dagger_\infty\to(\mathbb{A}^1)^\dagger$ 上的对数映射 $f^\dagger_\infty:C^\dagger_\infty\to X^\dagger$，其中 $C_\infty$ 上的对数结构由 $f^{-1}_\infty(\partial X)\subseteq C_\infty$ 诱导，而 $\mathcal O_\infty$ 上的对数结构是 $\mathbb{A}^1$ 上对数结构的拉回。这在 $\mathcal O_\infty$ 的闭点上的纤维 $C^\dagger_0$ 上诱导出一个对数结构，而假设是它与原来给定的 $\operatorname{Spec}k^\dagger$ 上的 $C^\dagger_0\to X^\dagger_0$ 重合。于是由例 3.26 中对数光滑曲线的描述不难看出，这蕴含 $C^\dagger_\infty\to\mathcal O^\dagger_\infty$ 是对数光滑的。但这样一来，由上述形变理论，$C^\dagger_k\cong C^\dagger_\infty\times_{\mathcal O^\dagger_\infty}\mathcal O^\dagger_k$ 对每个 $h$ 都被唯一确定，从而 $C^\dagger_\infty$ 唯一。
% [Tutor Note: p.169 末原文作“is uniquely determined for each $h$”，按上下文应为“for each $k$”（逐有限阶 $k$ 唯一确定，并与上一段的 $[f_k:C_k/O_k\to X,\mathbf{x}^k]$ 对应）。此处照原文作 $h$，正确形式见本注；见 PDF p.184 图像。]
% [Tutor Note: p.167 的精确列 $0\to\Theta_{C^\dagger_0/k^\dagger}\to f_0^*\Theta_{X^\dagger_0/k^\dagger}=M\otimes_\Z\mathcal O_C\to\mathcal N_{f_0}\to 0$ 中，靶项写作 $\Theta_{X^\dagger_0/k^\dagger}$（下标 $X^\dagger_0/k^\dagger$），而 p.166 的蛇引理图表与 (2.20) 处的同类列写作 $\Theta_{X^\dagger/(\mathbb{A}^1)^\dagger}$；这是同一对象在两种底（$k^\dagger$ 与 $(\mathbb{A}^1)^\dagger$）下的写法，照译不改，见 PDF p.182 图像。另 p.166 与 p.167 原文分别用 $H^0(C,\mathcal N_{f_0})$ 与 $H^0(C_0,\mathcal N_{f_0})$ 表示同一空间，亦照原文。]
\end{proof}

\subsection{证明的完成}
我们现在完成定理 4.4 的证明，即证明
\[
N^{0,\mathrm{trop}}_{\Delta,\Sigma}=N^{0,\mathrm{hol}}_{\Delta,\Sigma}.
\]
这里我们固定了点 $P_1,\dots,P_s\in M_\Q$，使得所有经过 $P_1,\dots,P_s$ 的热带有理曲线都是简单的。回忆我们已经重新标度了格 $M$，使每个 $P_i\in M$。此外，在选定了良好的分解 $\mathcal P$ 之后，我们可以再次重新标度格，以确保对每条亏格零热带曲线 $h:(\Gamma,x_1,\dots,x_s)\to M_\R$（其中 $h(x_i)=P_i$），$\widetilde\Gamma$ 的每条边的像的仿射长度都能被其权整除，从而定理 4.14 的假设总成立。
于是，照常选取截面 $\sigma_1,\dots,\sigma_s:\mathbb{A}^1\to X$，从而得到点 $q_i=\sigma_i(0)\in X_0$ 之后，由定理 4.14，我们恰好得到 $N^{0,\mathrm{trop}}_{\Delta,\Sigma}$ 条亏格零、满足 $f(x_i)=q_i$ 的环面横截带标记对数曲线
\[
f:(C^\dagger,x_1,\dots,x_s)\to X^\dagger_0 .
\]
用 $\mathcal M^{0,\mathrm{log}}_{\Delta,\Sigma}(\sigma_1,\dots,\sigma_s)$ 记这些环面横截带标记对数曲线的集合。
另一方面，令 $K=k((t))$。由包含 $k[t]\hookrightarrow K$ 给出一个典范映射 $\operatorname{Spec}K\to\mathbb{A}^1_k$，而纤维积 $X\times_{\mathbb{A}^1_k}\operatorname{Spec}K$ 同构于 $X_\Sigma\times_k K$。此外，每个截面 $\sigma_i$ 于是定义了 $X_\Sigma\times_k K$ 的一个 $K$-值点，我们也把它记作 $\sigma_i$。

% [原文 p.170]
我们希望证明：$\mathcal M^{0,\mathrm{log}}_{\Delta,\Sigma}(\sigma_1,\dots,\sigma_s)$ 与 $K$ 上亏格零、满足 $f(x_i)=\sigma_i$ 的环面横截曲线
\[
f:(C,x_1,\dots,x_s)\to X_\Sigma\times_k K
\]
的集合 $\mathcal M^{0,\mathrm{hol}}_{\Delta,\Sigma}(\sigma_1,\dots,\sigma_s)$ 之间存在一一对应。这将证明该定理。确实，由于 $K$ 是特征零的代数闭域，而答案应当与我们所使用的特征零代数闭域无关，$X_\Sigma\times_k K$ 中这类环面横截曲线的数目就是 $N^{0,\mathrm{hol}}_{\Delta,\Sigma}$。
为证明这一一一对应，从 $\mathcal M^{0,\mathrm{log}}_{\Delta,\Sigma}(\sigma_1,\dots,\sigma_s)$ 中的对数曲线
\[
[f_0:C_0/O_0\to X,\mathbf{x}^0]
\]
出发。由定理 4.27，这给出曲线
\[
[f_\infty:C_\infty/O_\infty\to X,\mathbf{x}^\infty].
\]
由包含 $k[[t]]\subseteq K$，我们有自然映射 $\operatorname{Spec}K\to\mathcal O_\infty$，从而有 $\operatorname{Spec}K$ 上的曲线 $C=C_\infty\times_{\mathcal O_\infty}\operatorname{Spec}K$，连同映射
\[
f:C\to X\times_{\mathbb{A}^1}\operatorname{Spec}K=X_\Sigma\times_k K .
\]
这就是所需的 $K$ 上的环面横截曲线。
反之，假设给定模空间 $\mathcal M^{0,\mathrm{hol}}_{\Delta,\Sigma}(\sigma_1,\dots,\sigma_s)$ 中的曲线 $f:C\to X_\Sigma\times_k K$。由于
\[
K=\bigcup_{d=1}^{\infty}k((t^{1/d})),
\]
而 $C$ 在 $K$ 上有限型，事实上存在某个 $d$ 使 $C$ 定义在域 $k((t^{1/d}))$ 上。这给出如 (4.15) 中的图表，其中 $C^*=C$，$L=k((t^{1/d}))$。然后应用定理 4.24，在把 $L$ 换成 $L=k((t^{1/de}))$（对某个 $e$）、作基变换 $\mathbb{A}^1\to\mathbb{A}^1$、并爆破 $X$ 之后，我们得到图表 (4.17)。注意爆破 $\widetilde X\to X$ 对应于 $\mathcal P$ 的一个细分。这于是通过把 (4.17) 中的映射 $f$ 限制到 $\operatorname{Spec}R$ 的闭点上的纤维，得到环面横截前对数曲线
\[
f_0:(C_0,x^0_1,\dots,x^0_s)\to\widetilde X_0 .
\]
它有相伴的热带曲线 $h:\Gamma\to M_\R$，必然经过 $P_1,\dots,P_s$。由于 $P_1,\dots,P_s$ 是取为一般位置的，$h$ 必是简单的，于是由定理 4.24 的 (2)，$f_0$ 是环面横截对数曲线。
但我们已经分类了所有相伴热带曲线经过 $P_1,\dots,P_s$ 的环面横截对数曲线，因此这条对数曲线正是对 $N^{0,\mathrm{trop}}_{\Delta,\Sigma}$ 有贡献的那些曲线之一。特别地，这条对数曲线给出一个已经定义在 $\operatorname{Spec}k[[t]]$ 上的唯一族。于是我们事实上得到 $de=1$，且不需要基变换 $\mathbb{A}^1_k\to\mathbb{A}^1_k$。特别地，所有定义在 $K$ 上的环面横截曲线事实上都定义在 $k((t))$ 上，而这些曲线与环面横截对数曲线之间有一一对应。
这就完成了 Mikhalkin 定理的证明。

% [原文 p.171]
\subsection{参考文献与延伸阅读}
本章的论述基于 Nishinou 与 Siebert 的论文 [86]。该文覆盖了任意维数下的类似结果。原始证明由 Mikhalkin 在 [80] 中给出，他此前在 [79] 中宣布了该结果。在 Mikhalkin 最初的宣布之后，Shustin 也在 [105] 中给出了该结果的一个独立证明。在 $\PP^2$ 的情形，Gathmann 与 Markwig 在 [31] 中给出了一个纯组合的论证，其办法是证明热带曲线计数满足 WDVV 方程。涉及 Mikhalkin 结果之其它推广的文献正在不断增多，这里无法一一列举。
% [Tutor Note: §4.7 的文献对应：Nishinou–Siebert [86]（任意维数的类似结果）、Mikhalkin [80]（原始证明，结果先在 [79] 宣布）、Shustin [105]（独立证明）、Gathmann–Markwig [31]（$\PP^2$ 情形的纯组合论证，证明热带曲线计数满足 WDVV 方程）。按 BRIEF 只保留文献编号，不逐条译文献。]
% [Tutor Note: §4.6 的枢纽等式 $N^{0,\mathrm{trop}}_{\Delta,\Sigma}=N^{0,\mathrm{hol}}_{\Delta,\Sigma}$（p.169）中，两个 $N$ 均为第 4 章记号（斜体 $N$），不是花体 $\mathcal N$（后者在本块中用于法层 $\mathcal N_{f_0}$）。另 p.169 原文把 $\mathcal M^{0,\mathrm{log}}_{\Delta,\Sigma}(\sigma_1,\dots,\sigma_s)$ 的元记为 $[f_0:C_0/O_0\to X,\mathbf{x}^0]$，与定理 4.27 的记号一致；p.170 末的 "$de=1$" 指基扩张次数 $d$ 与 $e$ 的乘积，照译。]
